\subsection{Ciphers for Asymmetric Encryption}\label{sec:asymmetric_cryptography:ciphers}

Below, we briefly present a short list of renowned or modern asymmetric encryption ciphers. Note that the use of asymmetric cryptography ciphers for digital signatures is covered in \Cref{sec:digital_signatures:ciphers}.

\subsubsection{Rivest-Shamir-Adleman}\label{sec:asymmetric_cryptography:ciphers:rsa}

\gls{rsa} \cite{rivest_method_1978} is an asymmetric cipher which relies on the hardness of the \gls{if} problem (\Cref{sec:trapdoor_functions:integer_factorization}). Being based on the \gls{if} problem, \gls{rsa} involves a modulus $n$ determined as the product of two prime numbers $p$ and $q$, a public exponent $e$, and a private exponent $d$. The size of the modulus $n$ in bits may vary and is usually one among 1024 (insecure), 2048, 3072, and 4096. A public \gls{rsa} key is a tuple $(n, e)$, while a private \gls{rsa} key is a tuple $(n, d)$.

\remarkBlock{On the value of the public exponent $e$}{The public exponent $e$ is usually set to 65537. Besides other mathematical reasons, the binary representation of 65537 has only two bits set to \texttt{1} --- i.e., 10000000000000001 --- and this allows for carrying out the exponentiation operation for encryption faster. Also, 65537 is a Fermat number (see \url{https://en.wikipedia.org/wiki/Fermat_number}).}

\gls{rsa} can be used for both encryption and digital signatures (\Cref{sec:digital_signatures}): RSAES-OAEP\footnote{RSAES-OAEP stands for \textit{RSA Encryption Scheme --- Optimal Asymmetric Encryption Padding} and is a secure way to encrypt small messages (or keys) with \gls{rsa} while protecting against chosen‐ciphertext attacks (\Cref{sec:security_of_cryptographic_algorithms})} is the best construction for encryption with padding, while RSASSA-PSS\footnote{RSASSA-PSS stands for \textit{RSA Signature Scheme with Appendix --- Probabilistic Signature Scheme} and is a secure way to generate digital signatures with \gls{rsa} by adding randomness (via a salt) and a provably secure padding construction to guard against forgery attacks (\Cref{sec:digital_signatures:security_properties})} is the recommended construction for digital signatures \cite{rfc8017}. In theory, \gls{rsa} can encrypt any message $m$ so that $m$ can be mapped to an integer between $0$ (included) and the modulus $n$ (excluded). Concretely, the message $m$ is padded into a block whose length in bits equals that of the modulus $n$. Depending on the padding scheme, the maximum practical message size may vary (e.g., OAEP allows messages to be up to the length of the modulus $n$ in bits - 528 --- for instance, 1520 bits if using a 2048-bit modulus). In reality, \gls{rsa} \gls{kem} (\Cref{sec:key_encapsulation_mechanisms}) --- which avoids padding altogether --- should be used for data encryption within hybrid cryptography (\Cref{sec:hybrid_cryptography}).

\warningBlock{Do not use \gls{rsa} unless you really have a good reason to}{The use of \gls{rsa} --- which is still popular in 2025 --- is now heavily discouraged, as more secure or efficient alternatives (e.g., \gls{ecc} --- see \Cref{sec:trapdoor_functions:elliptic_curve_cryptography}) are available. In fact, albeit involving fairly easy arithmetic operations, \gls{rsa} is very difficult to implement correctly due to concerns with prime numbers generation and choice of parameters (\Cref{sec:trapdoor_functions:integer_factorization}), and padding issues and side-channel attacks (\Cref{sec:security_of_cryptographic_algorithms}).}

\remarkBlock{\gls{rsa} and \gls{he}}{\gls{rsa} is multiplicative homomorphic. However, in practice the use of \gls{rsa} for \gls{he} is very limited due to the fact that padding invalidates \gls{rsa}'s homomorphism (\Cref{sec:advanced_cryptography:homomorphic_encryption}).}

\subsubsection{Paillier}\label{sec:asymmetric_cryptography:ciphers:paillier}
		
Paillier \cite{stern_public_1999} is an asymmetric cipher which relies on the hardness of the \gls{cr} problem (\Cref{sec:trapdoor_functions:composite_residuosity}). Being based on the \gls{cr} problem, Paillier involves a multiplicative cyclic group (based on integers), a modulus $n$ determined as the product of two prime numbers $p$ and $q$, and a generator $g$ (a common choice is $g = n + 1$). The size of the modulus $n$ in bits may vary and is usually one between 2048 and 3072. A public Paillier key is a tuple $(n, g)$, while a private Paillier key is a tuple $(\lambda, \mu)$ --- where $\lambda$ is the least common multiple between $p-1$ and $q-1$ and $\mu$ is $\bigl(L(g^\lambda \bmod n^2)\bigr)^{-1}\bmod n$ with $L(u) = \frac{u - 1}{n}$.

Paillier can be used for both encryption and digital signatures (\Cref{sec:digital_signatures}). Differently from \gls{rsa}, Paillier's encryption does not typically include padding. In theory, Paillier can encrypt any message $m$ so that $m$ can be mapped to an integer between $0$ (included) and the modulus $n$ (excluded). Concretely, the message $m$ is mixed with a random value $r$ picked randomly from the multiplicative cyclic group of nonzero elements $\mathbb{F}\textsuperscript{*}\textsubscript{n}$.

\remarkBlock{Paillier and \gls{he}}{Paillier is additive homomorphic, hence it can be used in use cases such as secure voting (tallies without revealing individual votes), privacy-preserving data aggregation (e.g., summing encrypted sensor readings) and threshold cryptography (e.g., cryptographic computations require collaboration among multiple parties --- see \Cref{sec:advanced_cryptography:homomorphic_encryption}.}

\subsubsection{ElGamal}\label{sec:asymmetric_cryptography:ciphers:elgamal}

ElGamal \cite{elgamal_public_1985} is an asymmetric cipher which relies on the hardness of the \gls{dl} problem (\Cref{sec:trapdoor_functions:discrete_logarithm}). Being based on the \gls{dl} problem, ElGamal involves a multiplicative cyclic group (either based on integers or elliptic curves --- see \Cref{sec:mathematical_underpinnings:elliptic_curves}), a prime modulus $p$, and a generator $g$. The size of the modulus $p$ in bits may vary and is usually one between 2048 and 3072 for integers and 256 or 384 for elliptic curves. A public ElGamal key is a value $h = g \bigcdot x$ --- where the operation $\bigcdot$ may be either exponentiation or point addition depending on whether the group is based on integers or elliptic curves, respectively --- while a private ElGamal key is the value $x$ (which belongs to the group).

ElGamal can be used for both encryption and (although seldom) also for digital signatures (\Cref{sec:digital_signatures}). Concerning encryption, ElGamal expects the encrypting party to choose an ephemeral (\Cref{sec:cryptography_basics:keys}) private key $k$ and derive a shared secret $s = h \bigcdot k$ --- shared with respect to the recipient's public key $h$ --- which in turn is used to encrypt the message $m$ as $c = m \bigcdot s$ (hence, $m$ must be an element of the group). The ciphertext and the ephemeral public key $(c, g \bigcdot k)$ are sent to the recipient for decryption. Differently from \gls{rsa}, ElGamal's encryption does not typically include padding.

\remarkBlock{ElGamal and \gls{he}}{ElGamal is multiplicative homomorphic, hence it can be used in use cases such as secure voting (tallies without revealing individual votes), privacy-preserving data aggregation (e.g., summing encrypted sensor readings) and threshold cryptography (e.g., cryptographic computations require collaboration among multiple parties --- see \Cref{sec:advanced_cryptography:homomorphic_encryption}.}

\remarkBlock{ElGamal and other asymmetric ciphers}{In a sense, ElGamal is a simpler version of \gls{ies} (\Cref{sec:hybrid_cryptography:ciphers:ies}) --- it is ``simpler'' because ElGamal's plaintexts must be (mapped to) an element of the group used in the underlying \gls{dl} problem, limiting flexibility, while in \gls{ies} this is not the case. Also, ElGamal is based on the same mathematics as the \gls{dh} key agreement (\Cref{sec:diffie_hellman}). However, differently from \gls{dh}, ElGamal immediately uses the shared secret to encrypt the message, while \gls{dh} uses the shared secret as one of the inputs to a \gls{kdf} (\Cref{sec:key_management:lifecycle:pre}) to derive a symmetric key (\Cref{sec:symmetric_cryptography}). Finally, ElGamal is the basis of the \gls{dsa} for digital signatures (\Cref{sec:digital_signatures:ciphers}) --- essentially, \gls{dsa} is a refined and standardized variant of ElGamal with well-defined parameters.}
